\section{Market value and book value}

It is natural to reason that depreciation gives the only change in capital value, once investment has ended, apart from market noise converging to zero, insofar as depreciation can be predicted, then, depreciated cost would show the norm (expectation) toward which current market value converges. It would follow that investors might anticipate the market by selling when prices are above depreciated cost, and buying when below.\footnote{As argued by James Tobin in 1968 and 1969. "Tobin's \(q\)" is the ratio of market to replacement value (RV). RV means depreciation cost under replacement cost accounting, where past investments and subsequent depreciation are re-expressed in the cheaper dollars of today. Investors were advised to invest when Tobin's \(q\) is less than one, and divest conversely.}

We accept that cost (investment) can realize equal market value, short of collective net investment, and that some firms neither innovate nor gain much from tech externalities, and that market noise accounts for the rest of change in value. But it appears to us that depreciated cost can find the norm described, even if depreciation were predicted perfectly, only insofar as did original cost. An argument that optimal accounting can give a sounder measure of worth than market value would seem  to need a supporting argument that price and worth were closer together at the time assets were acquired, and depreciation bases set at their cost, than over the depreciation period following. Garbage in, garbage out.

We do not question the value of depreciation accounting, or the importance of the depreciation concept in economic understanding. But we suggest that current market value gives a more reliable measure of capital, when both are known, than an assumption-dependent projection from past market value subject to the same criticisms





